\documentclass[11pt]{article}
\usepackage[a4paper,margin=2.4cm]{geometry}
\usepackage{newtxtext,newtxmath}
\usepackage{hyperref}
\pagestyle{empty}
\setlength{\parskip}{6pt}
\begin{document}

\noindent
To the Editors,\\
\emph{International Journal of Heat and Mass Transfer}

\vspace{10pt}
\noindent
Dear Editors,

\noindent
Please consider our manuscript, ``The Boussinesq--compressible gap in steady two-dimensional
buoyancy-driven enclosure convection is linear to leading order in the relative temperature
contrast: a factorised physics prior for surrogates'', for publication in the \emph{International Journal of
Heat and Mass Transfer} as a full-length article.

\noindent
Surrogates for thermal convection are trained overwhelmingly on Boussinesq data, yet the flows
engineers need are compressible. The two are close in integral quantities such as the Nusselt
number and differ in the local fields that surrogates are asked to reproduce. The manuscript
establishes what sets the size of that difference and turns the answer into a design rule.

\begin{itemize}\setlength{\itemsep}{2pt}
\item \textbf{One axis.} In the non-dimensionalisation used by the released paired data and by our
solver, the Froude number is pinned ($Fr^2=\varepsilon$ identically) and the Mach number enters at
$O(Ma^2)$; a 24-fold change in $Ma$ leaves the measured correction unchanged at the precision of
the solve. The
break-symmetry group is therefore one-dimensional, spanned by $\varepsilon=\Delta T/T_0$.
\item \textbf{A measured law.} Across 627 paired two-dimensional solves in the unit square at
$Pr=0.71$ spanning $\varepsilon=0.02$--$1.9$, the relative correction obeys
$R=A\varepsilon^{m}$ with $m=1.02$ in $u_x$, $u_y$ and $p$ (and the same exponent, with larger
scatter, in $\theta$), a vanishing intercept, and an $\varepsilon^2$ curvature that matches the
asymptotic expansion. In 104-sample checks in a 2:1
geometry and at $Pr=4.38$ the exponent is statistically unchanged; the amplitude reproduces on an
independently generated public data set at $\varepsilon=1.2$.
\item \textbf{A design rule.} Because the correction factorises as $D=\varepsilon P(X)$, a network
that predicts only $P$ has a $1.09\times$ lower channel-averaged error than the strongest
conditional baseline in extrapolation ($\theta$ and $p$ individually; the $u_x/u_y$ intervals
include zero), and $1.89\times$ better than an $\varepsilon$-blind network when $\varepsilon$ is
withheld, which makes
single-contrast public data reusable at other contrasts within one boundary-condition family, without
retraining.
\item \textbf{Wall heat transfer.} At $\varepsilon=1.2$ the integral Nusselt number changes by
$0.6\%$, while the local wall heat flux is redistributed by $4.25\%$; the integral agreement
therefore conceals a local error that a surrogate must be trained to resolve.
\end{itemize}

\noindent
IJHMT has published numerical methodology for non-Oberbeck--Boussinesq flows (Demou, Frantzis and
Grigoriadis, \emph{Int.\ J.\ Heat Mass Transfer} 125 (2018) 1156--1168), so the present study fits
the journal's scope: it quantifies that effect and turns it into a surrogate design rule.

\noindent
The main sweep is two-dimensional and steady; a preliminary twelve-sample three-dimensional check reproduces the
exponent, and a POD-ROM with at most four modes on the same samples lowers the mean $u_x$ error from 0.118 to
0.054 over four held-out contrasts. We state the remaining limits (one Rayleigh number, $res=16$,
$\varepsilon\le1.2$) explicitly. Every number in
the paper is regenerated from stored solver output by the scripts in the repository; the public
data set of the companion study is used unchanged, and our scripts and machine-readable results are
released with the manuscript; the raw solver output is available from the authors on request.

\noindent
The work is not under consideration elsewhere, and the authors declare no competing interests.

\vspace{12pt}
\noindent
Yours sincerely,\\
{[Corresponding author name]}\\
{[Affiliation]}\\
{[Email]}

\end{document}
